\section{Security}
\label{sec:security}

This section covers the security properties that follow from the architecture---value
conservation, consensus safety, and the bounded trust surface of the single-operator
venue---and the operational protections layered on top.

\subsection{Value conservation and consensus safety}

Two properties are structural rather than discretionary. \emph{Value conservation} follows
from atomic shared-memory import/export (Section~\ref{sec:cross-chain}): import and export
are applied in a single commit at block acceptance, so no path creates or destroys value,
and a block that is verified then rejected applies nothing. \emph{Consensus safety} follows
from the determinism discipline (Section~\ref{sec:determinism}): accepted state is a pure
function of the block, so all validators agree on the state root, and a proposer cannot
carry fabricated fills past a validator that re-derives them in \texttt{Verify}. Both are
exercised by tests---an end-to-end conservation test and same-block state-root equality
across fresh VMs---that fail on regression.

\subsection{Bounded trust surface of the single venue}

The D-Chain is a single-operator venue today. The residual trust this introduces is bounded
and named. A block the proposer builds and relays that then loses the consensus poll can
strand a D-Chain match with no settling counterpart; the blast radius is one taker's escrow,
and there is no supply inflation because conservation holds on every committed path. The
trustless rollup (Section~\ref{sec:rollup}) closes this surface by binding finality to a
verified validity proof rather than to proposer behavior, which is the security reason to
move from P3Q attestation to \texttt{starkfri} validity proofs as the venue decentralizes.

\subsection{Post-quantum verification}

The rollup verifiers are post-quantum. P3Q is FIPS-204 ML-DSA, and \texttt{starkfri} is a
STARK/FRI verifier over the Goldilocks field with a cSHAKE256 transcript---hash-based, with
no reliance on pairings or discrete-log hardness. Finality therefore does not rest on
assumptions broken by a quantum adversary, and the precompile envelope for \texttt{starkfri}
carries formal-verification artifacts for byte-level identity.

\subsection{Validator incentives and slashing}

Where the D-Chain decentralizes to multiple validators, the staking and slashing model
follows the Lux primary network: validators stake, earn a share of fees, and are slashed for
provable misbehavior. The relevant provable faults here are producing a block whose carried
fills or state root disagree with honest re-execution (a determinism violation), and
censorship of valid orders. Because the offense is detectable by deterministic
re-execution, slashing does not require a subjective judgment.

\subsection{Circuit breakers and oracle feeds}

Operational protections sit above the protocol. A volatility circuit breaker halts a market
when its price moves beyond a configured threshold within a window, resuming after a cooldown
or by governance action; the halt and resume are themselves consensus actions, so they apply
uniformly. Mark prices for liquidations and funding---where derivative markets are
enabled---are taken from decentralized oracle feeds with outlier rejection and on-chain
verification, rather than from a single source.
